% Bounded structural example from one of the eight unresolved HH cases.
% No main-manuscript edit, no HC search, no K2 single-edge test.
% Source phase6 SHA 7a90f79db489447b0bad3019217fd369a5cc25986288247981cbe7f3462acba3.
\paragraph{Common-triangle retention.}
If two selected cycles each retain two edges of the same common triangle,
and its three edges belong to neither pool, then they share at least
$2+2-3=1$ protected edge. This assertion only assumes the stated local
retention; it does not assert that arbitrary Hamilton cycles retain two
edges of every triangle.

For the interface $0112\ast\to0212\ast$ obtained by changing position
$1$ to colour $2$, the target normalizer is
$(s,\epsilon,p)=(3,-1,(0\ 2))$, with canonical target $0102\ast$.
Fix physical positions $0,1,2,4$ to $3,0,2,2$. At position $3$ both
neighbours and the excluded boundary colour are $2$, leaving colours
$0,1,3$. The three shared vertices are therefore
\[
 T=\{30202,30212,30232\},
\]
which induce a triangle of recolourings at position $3$.
The source cycle has the path $[30212,30202,30232]$ in $T$.
The restored target cycle has the path $[30202,30232,30212]$;
its canonical words are $[20230,30230,10230]$.
These two path identities follow directly from the printed cycles
for $0112\ast$ and $0102\ast$. Neither printed pool has an edge
in this triangle after restoration. The common-triangle rule thus gives
a protected common edge, without specifying one in an interface table.

% Only this one of the eight remaining cases has the tested K3 capacity.
% Seven others have no such selected two-of-three triangle support.
